\documentclass[a4paper,11pt]{article}
\usepackage[misc]{ifsym}
\usepackage{xcolor}

\usepackage[activate={true,nocompatibility}]{microtype}
\usepackage{microtype}
\usepackage[T1]{fontenc}
\usepackage{amsthm}
\newtheorem{lemma}{Lemma}[section]
%\usepackage{mlmodern}
%\usepackage{Baskervaldx}
\usepackage[xcharter]{newtx}
\usepackage{fontawesome5}


\renewcommand{\baselinestretch}{1.06}
\newcommand{\subt}{ <: }
\newcommand{\kmod}{\normalfont\textsf{\textbf{mod}}}
\newcommand{\ectx}[1]{\;@\;{\color{cyan} #1}}
\newcommand{\elab}[1]{{\color{gray}\dashrightarrow #1}}

%\usepackage{MnSymbol}
\usepackage{inconsolata}
\usepackage[scaled=0.92]{sourcesanspro}
\usepackage{xcolor}
\usepackage[utf8]{inputenc}
\usepackage[english]{babel}
\usepackage[a4paper, left=1.33in, right=1.33in]{geometry}
\usepackage{calc}
\usepackage{graphicx}

\usepackage{titlesec}
\titleformat{\subsubsection}[runin]
  {\normalfont\itshape}
  {\hspace*{\parindent}\normalfont\bfseries(\thesubsubsection)}
  {1.5\wordsep}
  {}
  []

\titlespacing*{\section}{0pt}{2\baselineskip}{1\baselineskip}
\titlespacing*{\subsubsection}
  {0pt}{6pt}{\wordsep}  % left, before, after spacing

\usepackage{mathpartir}
\usepackage{amsmath}

\usepackage{tikz}
\usepackage{wrapfig}

\usepackage[backend=bibtex, style=alphabetic]{biblatex}
\addbibresource{notes.bib}
\renewcommand*{\bibfont}{\small}

\author{Aghilas Y. Boussaa \texttt{<aghilas.boussaa@ens.fr>}}

\title{Notes on modal effect types}
\begin{document}
\sloppy
\maketitle

\section{Adapting \emph{Complete and Easy Bidirectional Typechecking for Higher-Rank Polymorphism}}
\subsection{Introduction}
\subsubsection{Brief overwiew of the plan.}
We remove explicit type application from \textsf{METL}. We want to adapt the
algorithm of Dunfield and Krishnaswami~\cite{dunfield2013complete} that allows
instantiation with monotypes. I will start by considering \emph{intuitionnistic}
monotypes (i.e.\ without modalities) to start, guarded monotypes may actually
work.

\subsubsection{Why monotypes?}
The paper mentions in Section~2 that it sticks to predicative polymorphism for
monotypes as the subtyping relation for general system F is undecidable.
However, it may still be decidable by allowing modalities. Allowing modalities
also means that we need to insert boxing and unboxing, when substituting a type
with one that has a top-level modality.

\subsubsection{Which subtyping relation?}
The original algorithm relies heavily on a subtyping relation, which we shall
adapt to \textsf{MET} types. To support unification from \textsf{METL} to
\textsf{MET}, the subtyping algorithm also has an elaboration part which gives a
function from $A$ to $B$ when $A$ is a subtype of $B$: $\Gamma\vdash A\textsf{
  <: }B\dashv\Delta{\color{gray} \dashrightarrow M}$.

Support for modalities in subtyping will come later. I start by only allowing
different modalities on the top-level. And using equality of modalities inside
types when unifying.

\subsection{Copying the original subtyping relation}
\subsubsection{Contexts.}
We merge the definitions of contexts given by both papers in
Figure~\ref{contexts}. We use the hole notation, meaning that: $\Gamma[\Delta]$
is a context of the form: $\Gamma_L,\Delta,\Gamma_R$.
\begin{figure}[h]
  \[
  \Gamma,\Delta\;::=\;\cdot\mid \Gamma,\alpha:K \mid \Gamma,\hat{\alpha}:K \mid
  \Gamma,\hat{\alpha}=\tau \mid \Gamma, x : A
  \mid \Gamma,\textsf{\scriptsize\faLock}_\mu
  \mid \Gamma, \blacktriangleright_{\hat{\alpha}}
  \]
  \caption{Syntax of contexts}\label{contexts}
\end{figure}


\subsubsection{Subtyping and instantiation.}
We have the same algorithmic subtyping and instantiation rules to which we add
the rule \textsc{Sub-Mod} to cross modalities when they are the same on both
sides; there is no rule for instantiation with modalities. And each rule has now
an elaboration part which is a (meta-theoretic) function, that maps terms to
terms.

When writing rules, we omit the kind of type variables when they are not
significant. In absence of kind annotation, a same variable should be considered
of having the same kind through the whole rule.

\begin{figure}[h]
  \centering
  \begin{mathpar}
    \inferrule[Sub-Var]{\\}{\Gamma[\alpha]\vdash\alpha\subt\alpha\dashv\Gamma[\alpha]
      {\color{gray} \dashrightarrow M\mapsto M}}
    
    \inferrule[Sub-ExVar]{\\}{\Gamma[\hat{\alpha}]\vdash\hat{\alpha}\subt\hat{\alpha}\dashv\Gamma[\hat{\alpha}]
      {\color{gray} \dashrightarrow M\mapsto M}}

    \inferrule[Sub-Unit]{\\}{\Gamma\vdash 1\subt1\dashv\Gamma
      {\color{gray} \dashrightarrow M\mapsto M}}

    \inferrule[Sub-Fun]{\Gamma\vdash B_1\subt A_1\dashv\Theta
      {\color{gray} \dashrightarrow f}\\
      \Theta{\vdash} [\Theta]A_2 {\subt} [\Theta]B_2\dashv\Delta
      {\color{gray} \dashrightarrow g}}{
      \Gamma\vdash A_1\to A_2\subt B_1\to B_2\dashv \Delta \elab{M\mapsto\lambda x^{B_1}.g (M (f(x)))}}

    \inferrule[Sub-ForallL]{\Gamma,\blacktriangleright_{\hat{\alpha}},{\hat{\alpha}:K}
      {\vdash}[{\hat{\alpha}}/\alpha]A\subt B\dashv\Delta,\blacktriangleright_{\hat{\alpha}},\Theta
      \elab{f}\\
      {\color{gray} \tau = [\Theta[\hat{\beta}=1/\hat{\beta}:K]]\hat{\alpha}}}
      {\Gamma\vdash\forall\alpha:K.A\subt B\dashv\Delta
       \elab{M \mapsto f(M \tau)}}

  \inferrule[Sub-ForallR]{\Gamma,\alpha:K\vdash A\subt B\dashv\Delta,\alpha:K,\Theta
    {\elab{f}}}
    {\Gamma\vdash A\subt \forall\alpha:K.B\dashv\Delta
      {\color{gray} \dashrightarrow M \mapsto \Lambda \alpha^K. f(M)}}

  \inferrule[Sub-InstantiateL]{\hat{\alpha}\not\in\textsf{FV}(A)\\
    {\Gamma}[\hat{\alpha}]\vdash\hat{\alpha}:\leqq A\dashv\Delta
    \elab{f}}
    {{\Gamma}[\hat{\alpha}]\vdash\hat{\alpha}\subt A\dashv\Delta
      \elab{f}}

  \inferrule[Sub-InstantiateR]{\hat{\alpha}\not\in\textsf{FV}(A)\\
    {\Gamma}[\hat{\alpha}]\vdash A\leqq:\hat{\alpha}\dashv\Delta
    \elab{f}}
    {{\Gamma}[\hat{\alpha}]\vdash A\subt \hat{\alpha}\dashv\Delta
      \elab{f}}

  \inferrule[Sub-Mod]{\Gamma\vdash G\subt G'\dashv\Delta\elab{f}}
    {\Gamma\vdash\overline{\mu} G\subt\overline{\mu} G'\dashv\Delta
     \elab{M\mapsto\textsf{\textbf{let mod}}_{\overline{\mu}} x = M \textsf{\textbf{ in }} \kmod_{\overline{\mu}} f(x)} }
  \end{mathpar}

  \caption{Algorithmic subtyping: $\Gamma\vdash A\subt B\dashv\Delta{\color{gray} \dashrightarrow M}$}
\end{figure}

The adaptation is straightforward expect for rule \textsc{Sub-ForallL} and its
type $\tau = [\Theta[\hat{\alpha}:K/\hat{\alpha}=1]]\hat{\alpha}$. The
substitution replaces all trailing existentials with unit types as their
definition is not important. For instance, when checking
$\forall\beta:\textsf{Any}.1<:1$, the rule inserts an existential variable whose
actual definition is useless. The type $\tau$ is an arbitrary solution.
\begin{figure}[h]
  \begin{mathpar}
    \inferrule[InstL-Solve]{\Gamma\vdash\tau:K}
      {\Gamma,\hat{\alpha}:K,\Gamma'\vdash\hat{\alpha}:\leqq\tau\dashv\Gamma,\hat{\alpha}=\tau,\Gamma'
      \elab{M\mapsto M }}

    \inferrule[InstL-Reach]{\\}
      {{\Gamma}[\hat{\alpha}:K][\hat{\beta}:K'] \vdash\hat{\alpha}:\leqq\hat{\beta}\dashv\Gamma[\hat{\alpha}:K\wedge K'][\hat{\beta}=\hat{\alpha}]
        \elab{M\mapsto M}}

    \inferrule[InstL-Arr]
      {{\Gamma}[{\hat{\alpha}}_2 :\textsf{Any}, {\hat{\alpha}}_1 :\textsf{Any},
                \hat{\alpha}={\hat{\alpha}}_1\to{\hat{\alpha}}_2]
        \vdash A_1\leqq: {\hat{\alpha}}_1 \dashv\Theta
        \elab{f}\\
        \Theta\vdash {\hat{\alpha}}_2:{\leqq}[\Theta]A_2\dashv\Delta
        \elab{g}}
      {{\Gamma}[{\hat{\alpha}:\textsf{Any}}]\vdash\hat{\alpha}:\leqq A_1 \to A_2\dashv\Delta
        \elab{M\mapsto\lambda x^{A_1}.g (M f(x))}}

    \inferrule[InstL-AllR]
      {{\Gamma}[\hat{\alpha}],\beta:K\vdash\hat{\alpha}:\leqq B\dashv\Delta,\beta:K,\Delta'
        \elab{f}}
      {{\Gamma}[\hat{\alpha}]\vdash\hat{\alpha}:\leqq \forall\beta:K.B\dashv\Delta
        \elab{M\mapsto\Lambda\beta^{K}.f(M)}}
  \end{mathpar}

\begin{mathpar}
    \inferrule[InstR-Solve]{\Gamma\vdash\tau:K}
      {\Gamma,\hat{\alpha}:K,\Gamma'\vdash\tau\leqq:\hat{\alpha}\dashv\Gamma,\hat{\alpha}=\tau,\Gamma'
      \elab{M\mapsto M}}

    \inferrule[InstR-Reach]{\\}
      {{\Gamma}[\hat{\alpha}:K][\hat{\beta}:K'] \vdash\hat{\beta}\leqq:\hat{\alpha}\dashv\Gamma[\hat{\alpha}:K\wedge K'][\hat{\beta}=\hat{\alpha}]
        \elab{M\mapsto M}}

    \inferrule[InstR-Arr]
      {{\Gamma}[{\hat{\alpha}}_2:\textsf{Any}, {\hat{\alpha}}_1 : \textsf{Any},
                 \hat{\alpha}={\hat{\alpha}}_1\to{\hat{\alpha}}_2]
        \vdash {\hat{\alpha}}_1:\leqq A_1 \dashv\Theta
        \elab{f}\\
        \Theta\vdash [\Theta]A_2{\leqq}:{\hat{\alpha}}_2\dashv\Delta
        \elab{g}}
      {{\Gamma}[{\hat{\alpha}}:\textsf{Any}]\vdash A_1 \to A_2\leqq:\hat{\alpha}\dashv\Delta
        \elab{M\mapsto\lambda x^{{\hat{\alpha}}_1}.g (M f(x))}}

    \inferrule[InstR-AllL]
      {{\Gamma}[{\hat{\alpha}}],\blacktriangleright_{\hat{\beta}},{\hat{\beta}:K}
      {\vdash}[{\hat{\beta}}/\beta]B\leqq: \hat{\alpha}\dashv\Delta,\blacktriangleright_{\hat{\beta}},\Delta'
      \elab{f}\\
      {\color{gray} \tau = [[\hat{\gamma}=1/\hat{\gamma}:K]\Delta']\hat{\beta}}}
      {{\Gamma}[\hat{\alpha}]\vdash\forall\beta:K.B\leqq:\hat{\alpha}\dashv\Delta
        \elab{M\mapsto f(M\tau)}}
  \end{mathpar}

  \caption{Instantiation}
\end{figure}

The following two lemmas guarantee the soundness of elaboration.
\begin{lemma}
  If $\Gamma\vdash A\subt B\dashv\Delta{\color{gray} \dashrightarrow f}$, and
  $\Gamma\vdash M:A\ectx{E}$, then $\Delta\vdash f(M):[\Delta]B$.
\end{lemma}
\begin{lemma}
  If $\Gamma\vdash A\subt B\dashv\Delta{\color{gray} \dashrightarrow f}$, then
  $f$ maps values to values.
\end{lemma}

\subsubsection{New inference rules.}
We take the rules of \textsf{METL} (Appendix~D of~\cite{tang2025modal}) as a
starting point, remove the rules of both application and type application and
try to add the new inference rule for application --- I will try to add
inference for abstraction later. We also give a new rule for switching from
checking to inference mode.

The rule \textsc{B-Switch} combines the switching from both algorithms. We add a
new rule for application (\textsc{App-Mod}) that deals with modalities in the
same way the original \textsf{METL} \textsc{B-App} rule did. This formulation
allows interleaving modalities and quantifiers.
\begin{figure}[h]
  \begin{mathpar}
    \inferrule[B-Switch]
      {\Gamma\vdash M\Rightarrow\overline{\mu}G'\ectx{E}
        {\color{gray} \dashrightarrow M'}\\
       (G',\overline{\mu})\Rightarrow\overline{\nu}\ectx{E}\\
        \Gamma\vdash G' <: G \dashv \Delta {\color{gray} \dashrightarrow f}
        \\
        \text{$M$ is not a value or $\overline{\nu}$ is empty}
      }
      {\Gamma\vdash M \Leftarrow\overline{\nu}G \dashv\Delta\ectx{E}
       \elab{\textsf{\textbf{let mod}}_{\overline{\mu}} x = M'
        \textsf{ in } \kmod_{\overline{\nu}} f(x)}}

    \inferrule[B-App]
      {\Gamma\vdash M\Rightarrow A\dashv\Delta\ectx{E}
       {\color{gray} \dashrightarrow M'}\\
       \Delta\vdash ({\color{gray}M':\ } A) \bullet N \Rrightarrow
       C\dashv\Theta\ectx{E}
       {\color{gray} \dashrightarrow M''}}
      {\Gamma \vdash M N:C\dashv\Theta\ectx{E}
        {\color{gray} \dashrightarrow M''}}

    \inferrule[App-Fun]
      {\Gamma\vdash N\Leftarrow A\dashv\Delta\ectx{E}
       {\color{gray} \dashrightarrow N'}}
      {\Gamma\vdash ({\color{gray}M:\ }A\to C) \bullet N \Rrightarrow C\dashv\Delta\ectx{E}
       \elab{M N'}}

    \inferrule[App-ExVar]
      {{\Gamma}[{\hat{\alpha}}_2:\textsf{Any}, {\hat{\alpha}}_1 : \textsf{Any},
                 \hat{\alpha}={\hat{\alpha}}_1\to{\hat{\alpha}}_2]
        \vdash N {\Leftarrow} {\hat{\alpha}}_1\dashv\Delta
        \ectx{E}\elab{N'}}
      {{\Gamma}[\hat{\alpha}:\textsf{Any}]\vdash
        ({\color{gray}M:\ } \hat{\alpha})\bullet N{\Rrightarrow}
        {\hat{\alpha}}_2\dashv\Delta\ectx{E}\elab{M N'}}

    \inferrule[App-All]
      {\Gamma,\hat{\alpha}:K\vdash ({\color{gray}M\hat{\alpha}:\ } [\hat{\alpha}/\alpha]A)
        \bullet N\Rrightarrow C\dashv\Delta\ectx{E}\elab{M'}}
      {
        \Gamma\vdash ({\color{gray}M:\ } \forall\alpha:K.A) \bullet N \Rrightarrow C\dashv\Delta\ectx{E}
        \elab{M'}
      }

    \inferrule[App-Mod]
      {{\overline{\mu}}_E\Rightarrow\langle\rangle_E\\
        {\color{gray} M'=\textsf{unmod}(M;\overline{\mu})}\\
        \Gamma\vdash ({\color{gray}M':\ } G) \bullet N \Rrightarrow C\dashv\Delta\ectx{E}
        \elab{M''}}
      {{\Gamma}\vdash
        ({\color{gray} M :\ } \overline{\mu}G)\bullet N{\Rrightarrow} C
        \dashv\Delta\ectx{E}\elab{M''}}
  \end{mathpar}
  \caption{Algorithmic typing}
\end{figure}

We omitted the inference rules that relying on the join operation, but we modify
its definition by redifining it on types. We state the rules for joining types
in Figure~\ref{join}. The join operation is now presented as a judgement that
outputs a new context and a coercion function. This new operation is then used
by the elaboration function \textsf{join}, which should insert coercion
functions.

The rule \textsc{Join-Abs} is correct because of the following lemma that
guarantees that $[\Delta]A$ is pure if and only if $[\Delta]B$.
\begin{lemma}
  If $\Gamma\vdash A<:B \dashv\Delta$ and
  $\Delta{\vdash}[\Delta]B:{\normalfont\textsf{Abs}}$, then
  $\Delta{\vdash}[\Delta]A:{\normalfont\textsf{Abs}}$.
\end{lemma}
\begin{figure}[h]
  \begin{mathpar}
    \inferrule[]
      {\Gamma\vdash B<: A\dashv\Delta\elab{f}\\
       \Delta{\vdash} [\Delta]A<: [\Delta]B\dashv\Theta\elab{g}}
      {\Gamma\vdash A:=B\dashv\Theta\elab{(f, g)}}

    \inferrule[Join-Abs]
      {\Gamma\vdash G:=G'\dashv\Theta\elab{(f, g)}\\
       \Delta{\vdash}[\Delta]G:\textsf{Abs}}
      {\Gamma\vdash\overline{\mu}G\vee\overline{\nu}G' = [\Delta]G\dashv\Delta\ectx{F}
       \elab{(f, g)}}

    \inferrule[Join-Mod]
      {\Gamma\vdash G:=G'\dashv\Theta\elab{(f, g)}\\
       \xi=\overline{\mu}\vee_{\Gamma,F}\overline{\nu}\\
       \color{gray}
       g'(M)=\textbf{\textsf{ let mod}}_{\overline{\mu}} x = M \textbf{\textsf{ in }} \kmod_\xi g(x)\\
       \color{gray}
       f'(M)=\textbf{\textsf{ let mod}}_{\overline{\nu}} x = M \textbf{\textsf{ in }} \kmod_\xi f(x)\\
      }
      {\Gamma\vdash\overline{\mu}G\vee\overline{\nu}G' = [\Delta]\xi G\dashv\Delta\ectx{F}
       \elab{(f',g')}}
  \end{mathpar}
  \caption{Join operation on types}\label{join}
\end{figure}

\subsection{Instanciation with modalities}
\subsubsection{Subtyping modalities}
We start by modifying the subtyping relation to allow subtyping of types with
modalities, generalising the rule \textsc{B-Switch}. The subtyping relation now
depends on an environment context: $\Gamma\vdash A\subt B \dashv \Delta \, @ \,E
{\color{gray} \dashrightarrow f}$. Rules apart from \textsc{Sub-Mod} keep the
same context.
\[
\inferrule[Sub-Mod]{\\}{\Gamma\vdash\overline{\mu}G\subt\overline{\nu}G'\dashv\Delta\ectx{E}}
\]

\subsubsection{Problems.}
When unifying a flexible type variable $\hat{\alpha}$ against a
non-intuitionistic monotype $\tau$, producing the equation $\hat{\alpha}=\tau$
is not complete. For instance, we may check $\hat{\alpha}$ first against the
type $[\textsf{Gen int}](1\to1)$, and then against $<\textsf{Gen int}>(1\to1)$.

A complete algorithm may not exist, but we may do better than plain assignation
by recording inequalities modulo effect context in the context:
$\hat{\alpha}\leq \tau\ectx{E}$. Elaboration becomes more involved as there are
now holes with functions that can only be filled when the type variable gets a
definition.

\end{document}
